\section{Engineering spin control and exchange interactions}
\label{sec:spin-control-exchange}

An encoding specifies which states store quantum information.
To manipulate that information, we need a Hamiltonian that changes
relative phases or populations within the chosen spin space.
A static magnetic field splits the two levels of an individual
spin and makes their relative phase evolve. A transverse drive
mixes those levels, while tunneling between dots produces an
exchange interaction that couples different spins. These physical
terms supply the controls whose action on each encoding is worked
out in Section~\ref{sec:quantum-gates}.

The spin description assumes that charge, orbital, and valley
excitations can be excluded as independent computational states.
This restriction must remain valid at the operating field and
under the applied pulses, even when virtual excitations contribute
to the effective spin Hamiltonian~\cite{Burkard2023}.

\subsection{Zeeman splitting and the two-level approximation}
\label{sec:platform-encoding}

Consider a singly occupied dot with the electron in the spatial
state $\ket{\chi_0}$. A static magnetic field along $+z$ lifts the
degeneracy of its two spin states. For a positive isotropic
effective $g$ factor, the Zeeman Hamiltonian in the physical-spin
basis $(\ket{\uparrow},\ket{\downarrow})$ is
\begin{equation}
  H_Z=\frac{E_Z}{2}\sigma_z,
  \qquad E_Z=g\mu_B B>0.
  \label{eq:zeeman}
\end{equation}
Here $\mu_B$ is the Bohr magneton and $B$ the field magnitude.
With the logical basis in Eq.~\eqref{eq:qubit-basis},
$\ket0$ is the lower level and $\ket1$ the upper level.
In this basis, $H_Z|_{\Hilb_{\mathrm q}}=-E_ZZ/2$.
The minus sign follows from assigning the lower spin level to
$\ket0$. This diagonal Hamiltonian changes the relative phase of
a spin superposition while leaving the two populations unchanged.

The two-level approximation requires the electron to remain in the
retained orbital--valley state and charge sector. Let
$\Delta_{\mathrm{orb}}$ and $\Delta_{\mathrm{val}}$ denote the
relevant gaps to excluded orbital and valley states, measured at the
operating field (Appendix~\ref{app:dot-spectroscopy}). The condition
$k_BT\ll\min(\Delta_{\mathrm{orb}},\Delta_{\mathrm{val}})$
suppresses their thermal occupation. If electron exchange with a
reservoir remains possible, every accessible charge-transfer process
out of either retained spin state must have a positive energy cost
large compared with $k_BT$. The costs $\Delta_\pm$ in
Eq.~\eqref{eq:dot-charge-transfer-costs} refer to the ground state
of the occupied charge sector, with $E_N$ evaluated including the
Zeeman energies at the operating field. Starting in its upper Zeeman state
reduces a transfer cost by $E_Z$ for the same final state.
Alternatively, reservoir exchange can be made negligible on the
operation timescale by raising its tunnel barrier. Low temperature
alone does not prevent
control-induced leakage: a pulse that couples to an excluded state
can populate it if its frequency approaches the corresponding
transition. The coupling matrix elements, excitation gaps, and
pulse spectrum must therefore be checked together~\cite{Burkard2023}.

Both Zeeman states belong to the retained qubit space. Keeping the
electron in that space does not require it to occupy only the
lower spin level, and does not by itself require $k_BT\ll E_Z$.
That stronger spin-polarization condition is useful when thermal
relaxation is used to prepare $\ket0$. During coherent operations,
the intended state may instead be any superposition of the two
levels. Spin relaxation and loss of phase coherence constrain how
long that state can be retained even when charge, orbital, and
valley leakage remain negligible.

\subsection{Transverse spin control}
\label{sec:transverse-spin-control}

Our aim is to control both the relative phase and the populations
of the two spin states. A suitable Hamiltonian is
\begin{equation}
  H_{\mathrm{target}}(t)=a(t)\sigma_z+b(t)\sigma_x,
  \label{eq:single-spin-control-target}
\end{equation}
where $a,b$ are real energy coefficients.
The $\sigma_z$ term changes the relative phase of $\ket{\uparrow}$
and $\ket{\downarrow}$, while $\sigma_x$ mixes them.
A static Zeeman field already supplies the longitudinal term.
We need to generate a useful transverse term, either directly in
the laboratory or in an effective rotating frame.

A constant transverse coupling is the simplest candidate. To see
its limitation, keep the Zeeman splitting fixed and switch on
a term $\hbar\Omega_x\sigma_x/2$ at $t=0$:
\begin{equation}
  H_{\mathrm{static}}=\frac{\hbar}{2}
     (\omega_Z\sigma_z+\Omega_x\sigma_x),\qquad
  \omega_Z=\frac{E_Z}{\hbar}.
  \label{eq:static-transverse-hamiltonian}
\end{equation}
Since $(\omega_Z\sigma_z+\Omega_x\sigma_x)^2
=(\omega_Z^2+\Omega_x^2)I$, a spin initially in $\ket{\downarrow}$
has transition probability
\begin{equation}
  P^{\mathrm{static}}_{\downarrow\to\uparrow}(t)
   =\frac{\Omega_x^2}{\omega_Z^2+\Omega_x^2}
      \sin^2\!\left(\frac{t}{2}\sqrt{\omega_Z^2+\Omega_x^2}\right),
  \qquad
  P_{\max}^{\mathrm{static}}
   =\frac{\Omega_x^2}{\omega_Z^2+\Omega_x^2}.
  \label{eq:static-transverse-probability}
\end{equation}
For $|\Omega_x|\ll\omega_Z$, the maximum is approximately
$(\Omega_x/\omega_Z)^2\ll1$, regardless of pulse duration. The spin rotates about $(\Omega_x,0,\omega_Z)$,
an axis close to $z$, so it cannot reach the opposite Zeeman state.
With zero longitudinal splitting, a full flip would take
$\pi/|\Omega_x|$; the bound applies to a single constant pulse
with fixed $\omega_Z$.

Resonant driving avoids this limitation by making the transverse
coupling oscillate near the Zeeman frequency.

\paragraph{Resonant driving and Rabi oscillations.}
Keep the Zeeman splitting fixed and take an oscillating transverse
term with real drive coefficient $\Omega$, angular frequency
$\omega$ and phase $\phi$ (Appendix~\ref{app:transverse-drive-origins}):
\begin{equation}
  H(t)=\frac{\hbar\omega_Z}{2}\sigma_z
       +\hbar\Omega\cos(\omega t+\phi)\sigma_x,
  \qquad \omega_Z=\frac{E_Z}{\hbar}.
  \label{eq:esr-lab}
\end{equation}
The drive couples the Zeeman states because
$\bra{\uparrow}\sigma_x\ket{\downarrow}=1$.
In a frame rotating at the drive frequency,
$\ket{\psi(t)}=e^{-i\omega t\sigma_z/2}\ket{\psi_{\mathrm r}(t)}$,
define the detuning $\delta=\omega_Z-\omega$.
For $|\delta|,|\Omega|\ll\omega_Z$, the rotating-wave approximation
(RWA) gives the effective Hamiltonian and evolution
(Appendix~\ref{app:spin-driving})
\begin{equation}
  \begin{aligned}
    H_{\mathrm{RWA}}&=\frac{\hbar}{2}
      \left[\delta\sigma_z
        +\Omega(\cos\phi\,\sigma_x+\sin\phi\,\sigma_y)\right],\\
    U_{\mathrm r}(t)&=\cos\frac{\Omega_{\mathrm{eff}}t}{2}\,I
      -i\sin\frac{\Omega_{\mathrm{eff}}t}{2}\,
       \frac{\delta\sigma_z
         +\Omega(\cos\phi\,\sigma_x+\sin\phi\,\sigma_y)}
            {\Omega_{\mathrm{eff}}},
  \end{aligned}
  \label{eq:esr-propagator}
\end{equation}
where $\Omega_{\mathrm{eff}}=\sqrt{\delta^2+\Omega^2}$ is the rotation
frequency. The phase sets the transverse direction of the rotation
axis $(\Omega\cos\phi,\Omega\sin\phi,\delta)$. For $\phi=0$,
the target coefficients are $a=\hbar\delta/2$ and $b=\hbar\Omega/2$.
To find the pulse duration for a spin flip, start in
$\ket{\downarrow}$ and evaluate the up-spin probability:
\begin{equation}
  P_{\downarrow\to\uparrow}(t)
    =\frac{\Omega^2}{\delta^2+\Omega^2}
       \sin^2\!\left(\frac{t}{2}\sqrt{\delta^2+\Omega^2}\right).
  \label{eq:rabi-probability}
\end{equation}
This population-transfer probability is independent of $\phi$.
Tune the drive frequency to the measured spin resonance to set
$\delta=0$ (Appendix~\ref{sec:measurement-spin-spectroscopy}).
The measured rotation rate at the chosen drive amplitude then sets
the spin-flip duration $t_\pi=\pi/|\Omega|$, even for
$|\Omega|\ll\omega_Z$.
Detuning tilts the effective rotation axis and reduces the maximum
inversion. High-fidelity inversion therefore requires
$|\delta|\ll|\Omega|$ and a pulse shorter than the coherence time.

\subsection{Exchange from virtual tunneling}
\label{sec:exchange-virtual-tunneling}

For electron spins confined in neighboring quantum dots, tunneling
through the interdot barrier provides an exchange interaction whose
strength can be controlled electrically. Together with the single-spin
controls described above, this interaction supplies the means to
implement two-spin gates. In the regime of interest, the electron
pair remains predominantly in the $(1,1)$ configuration: double occupation is energetically
costly, but its weak admixture can still affect the spin dynamics.

\paragraph{Virtual tunneling and exchange.}
Consider a double dot isolated from its reservoirs, so that its total
electron number remains fixed at two. Tunneling connects the $(1,1)$
configuration to $(2,0)$ and $(0,2)$ by transferring an electron onto
the occupied neighboring dot. To describe these processes, we retain
one orbital--valley state on each dot, with both spin states, and
assume that tunneling conserves spin. The Hubbard approximation
developed in Appendix~\ref{sec:device-hubbard-reduction} gives
\begin{equation}
  \begin{aligned}
    H_{\mathrm{Hub}}
      &=\sum_{i=1}^2\epsilon_i n_i
        +U\sum_{i=1}^2 n_{i\uparrow}n_{i\downarrow}+Vn_1n_2\\
      &\quad-t_c\sum_{\sigma=\uparrow,\downarrow}
        (c_{1\sigma}^\dagger c_{2\sigma}
         +c_{2\sigma}^\dagger c_{1\sigma}).
  \end{aligned}
  \label{eq:two-site-hubbard}
\end{equation}
Here $c_{i\sigma}$ annihilates an electron with spin $\sigma$ on
dot $i$. The occupation operators are
$n_{i\sigma}=c_{i\sigma}^\dagger c_{i\sigma}$ and
$n_i=n_{i\uparrow}+n_{i\downarrow}$.
Each retained orbital has energy $\epsilon_i$, while the Coulomb
repulsion depends on whether the electrons occupy the same dot or
different dots. These two cases contribute the energies $U$ and $V$,
respectively; for simplicity, we take the same $U$ on both dots.
The final term transfers an electron between the dots with real
tunnel amplitude $t_c$. This model neglects direct Coulomb exchange,
and we leave out the local Zeeman terms when calculating the exchange
contribution.

The energetic cost of a tunneling process follows from both the
orbital energies and the repulsion. Moving an electron onto the
occupied dot replaces the interdot repulsion $V$ with the on-site
repulsion $U$. It also changes the orbital energy by
$+\varepsilon$ or $-\varepsilon$, depending on the direction of
transfer, where $\varepsilon=\epsilon_1-\epsilon_2$.
The two double-occupation costs, measured from the energy of $(1,1)$
in the absence of tunneling, are therefore
\begin{equation}
  \Delta_{20}=U-V+\varepsilon,\qquad
  \Delta_{02}=U-V-\varepsilon.
  \label{eq:virtual-charge-gaps}
\end{equation}
We choose an operating point at which both costs are positive and
much larger than $|t_c|$. Tunneling then mixes only small components
of the doubly occupied states into the low-energy states associated
with $(1,1)$. Provided thermal and drive-induced charge excitation
are also negligible, the doubly occupied states can be eliminated
from the low-energy description while retaining the energy shifts
that their virtual mixing produces.

The spin dependence of these shifts follows from fermionic
antisymmetry. Within the retained orbital--valley state, a doubly
occupied dot can support only a spin singlet. The $(1,1)$ singlet
can therefore mix with both higher-energy charge states, whereas
the triplets have no such tunneling channel. This mixing lowers
the singlet energy and leaves the triplet energy unchanged.
Eliminating the two charge states to second order in $t_c$ gives
the singlet--triplet energy difference
(Appendix~\ref{sec:device-low-energy-reduction})
\begin{equation}
  J\equiv E_T-E_S
     =2t_c^2\left(\frac{1}{\Delta_{20}}+\frac{1}{\Delta_{02}}\right).
  \label{eq:superexchange}
\end{equation}
This energy difference is the exchange coupling that enters the
spin Hamiltonian. After the higher-energy charge states have been
eliminated, the interaction within the $(1,1)$ spin space takes
the Heisenberg form~\cite{Loss1998}
\begin{equation}
  H_{\mathrm{ex}}=J\,\mathbf{s}_1\cdot\mathbf{s}_2.
  \label{eq:exchange-spin}
\end{equation}
A common energy shift has been omitted. The charge degrees of freedom
thus enter the spin dynamics through the coefficient $J$, whose value
is set by the tunnel amplitude and the two double-occupation costs.

\paragraph{Tuning exchange with gate voltages.}
Gate voltages control this interaction by changing the interdot
tunneling and the relative orbital energies. Substituting the charge
costs into the expression for $J$ makes these dependencies explicit:
\begin{equation}
  J(\varepsilon)=\frac{4t_c^2(U-V)}{(U-V)^2-\varepsilon^2},\qquad
  J(0)=\frac{4t_c^2}{U-V}.
  \label{eq:exchange-detuning}
\end{equation}
Adjusting the interdot barrier changes the tunnel amplitude $t_c$,
while shifting the orbital energy of one dot relative to the other
changes the detuning $\varepsilon$. The tunnel amplitude at the
chosen barrier setting must be measured
(Appendix~\ref{sec:measurement-tunnel-spectroscopy}); the exchange
strength used to set a spin pulse is calibrated at the corresponding
compensated voltages (Appendix~\ref{sec:calibration-two-spin-exchange}).
In either case, the resulting
exchange pulse must remain within the regime used to obtain the
spin Hamiltonian. If a double-occupation cost becomes comparable
to $|t_c|$, the associated charge state mixes strongly with $(1,1)$
and must be included explicitly in the dynamics. The apparent
divergence of the expression for $J$ near a charge degeneracy is
therefore outside its range of validity.

In a device, changing the barrier voltage generally changes the
tunnel amplitude and shifts the dot levels at the same time. A pulse intended to tune
exchange can consequently move the addition thresholds that define
the chosen charge operating point. These shifts can be compensated
by adjusting the two plunger voltages together with the barrier
voltage.

For a small barrier-voltage change $\delta V_{\mathrm B}$, let
$\delta\mathbf V_{\mathrm P}$ denote the two plunger adjustments
and $\delta\boldsymbol\mu$ the resulting shifts of the two addition
thresholds. The matrix $G_{\mathrm P}$ contains the measured
derivatives of these thresholds with respect to the two plunger
voltages at the operating point, while the vector $\mathbf g_{\mathrm B}$
contains their measured derivatives with respect to the barrier
voltage. With independent plunger responses,
$G_{\mathrm P}$ is invertible, and the voltage response gives the
compensation required to keep both thresholds fixed
(Appendix~\ref{app:virtual-gates}):
\begin{equation}
 \delta\boldsymbol\mu
    =G_{\mathrm P}\delta\mathbf V_{\mathrm P}
          +\mathbf g_{\mathrm B}\delta V_{\mathrm B},
 \qquad
 \left.\delta\mathbf V_{\mathrm P}\right|_{\delta\boldsymbol\mu=0}
    =-G_{\mathrm P}^{-1}\mathbf g_{\mathrm B}\delta V_{\mathrm B}.
 \label{eq:platform-compensated-voltage}
\end{equation}
Coordinated voltage combinations chosen to control specified
quantities are called \emph{virtual gates}. This combination cancels
the changes in the two chosen thresholds to first order around the
operating point, while allowing the tunnel amplitude and other
Hamiltonian parameters to vary.

